\documentclass[a5paper,14pt]{report}
\usepackage{amsmath}
\usepackage{amssymb}
\usepackage[english, russian]{babel}
\usepackage[cp1251]{inputenc}
\usepackage[dvips]{graphicx}
\graphicspath{{images/}}
\usepackage{indentfirst}
\usepackage{multicol}
\usepackage{hhline}
\usepackage{multirow}
\usepackage{textcomp}
\usepackage{wrapfig}
\usepackage{color,colortbl}
\definecolor{Gray1}{gray}{0.7}
\definecolor{Gray2}{gray}{0.6}
\definecolor{Gray3}{gray}{0.5}
\definecolor{Gray4}{gray}{0.4}
\usepackage{geometry}
\geometry{left=1.5cm}
\geometry{right=1.5cm}
\geometry{top=1cm}
\geometry{bottom=2cm}

\begin{document}
\pagestyle{empty}
\begin{flushright}
\includegraphics[width=5cm]{ITMO_log}
\end{flushright}
\textbf{Миссия университета} -- генерация передовых знаний, внедрение инновационных разработок и подготовка элитных кадров, способных действовать в условиях быстро меняющегося мира и обеспечивать опережающее развитие науки, технологий и других областей для содействия решению актуальных задач. 
\\\line(1,0) {350}
\begin{center}
\textbf{КАФЕДРА ВЫЧИСЛИТЕЛЬНОЙ ТЕХНИКИ}
\end{center}
Кафедра вычислительной техники Университета ИТМО более 70 лет ведёт подготовку высококвалифицированных специалистов в области информатики и вычислительной техники. За годы своего существования на кафедре подготовлено более 5,5 тысяч высококвалифицированных специалистов, подготовлено более 240 кандидатов наук и 37 докторов наук. Преподавателями и научными сотрудниками кафедры изданы десятки монографий, учебников и учебно-методических пособий. Кафедра имеет высококвалифицированный состав преподавателей, среди которых 7 профессоров и 14 доцентов, обучающих более 500 студентов и аспирантов.
\\\\Более 15 лет кафедра ведёт подготовку по двухуровневой схеме: бакалавриат-магистратура. За эти годы подготовлено около полутора тысяч бакалавров и магистров по направлению «Информатика и вычислительная техника». В 2011 году на кафедре началась подготовка бакалавров и магистров по новому направлению «Программная инженерия». 
\\\\Ведущими преподавателями кафедры разработаны 7 магистерских программ, охватывающих наиболее актуальные направления вычислительной техники, такие как вычислительные системы и компьютерные сети, информационная безопасность и защита информации, встроенные вычислительные системы и системы на кристалле, программные комплексы и информационные интеллектуальные системы. Магистранты кафедры вычислительной техники активно участвуют в научно-исследовательских разработках кафедры, выступают с докладами на конференциях и семинарах различных уровней, публикуют свои результаты в научно-технических журналах и сборниках.
\end{document}